% 原文PDF405末尾supersymmetryからPDF410末尾zero-slope limitまで。
\gloss{超対称性（supersymmetry）}
{電荷がスピノルとして変換する対称性であり、フェルミオンとボソンの質量および結合を関係付ける。}

\gloss{超対称性の破れのスケール（supersymmetry breaking scale）}
{超対称性を破る期待値の質量スケール。超対称パートナーのスケールは、超対称性の破れのスケールに、標準模型の場と超対称性を破る場との結合の強さを掛けたものである。}

\gloss{T 双対性（T-duality）}
{通常はトーラスコンパクト化された弦理論における双対性で、半径に依存するスケール変換を除いて結合定数を不変に保ち、したがって弦の摂動論の各次数で成り立つ。最も顕著なのは $R\to\alpha'/R$ の双対性であり、巻き付き状態と Kaluza--Klein 状態を入れ替えることで、大きなトーラスと小さなトーラスにコンパクト化された弦理論を関係付ける。より一般には、反対称テンソル背景のシフトと、時空の大きな座標変換も含む。S 双対性および U 双対性と対照せよ。}

\gloss{'t Hooft--Polyakov モノポール（'t Hooft--Polyakov monopole）}
{磁荷をもつ古典解であり、単純群が $U(1)$ 因子をもつ群へ自発的に破れるときには必ず存在する。}

\gloss{タキオン（tachyon）}
{質量の二乗が負の粒子。ほとんどの場合はスカラーであり、真空の不安定性を意味する。}

\gloss{タドポール（tadpole）}
{量子効果によって誘起される、真空から単一粒子を生成する振幅。}

\gloss{標的空間（target space）}
{関数が値をとる空間。通常は弦の世界面上の非線形シグマ模型について用いられ、その場合、標的空間は時空そのものである。}

\gloss{Teichmüller パラメーター（Teichmüller parameters）}
{Riemann 面の複素構造のモジュライ。厳密には、Teichmüller 空間の点はモジュラー群の作用で同一視されない。}

\gloss{張力のない弦理論（tensionless string theory）}
{張力のない弦を含む、相互作用する理論。そのような弦は、縮みつつあるサイクルに $p-1$ 個の方向を巻き付けた $p$ ブレーンとして、高次元ブレーンが一致するときに、それらの間へ一つの方向を伸ばしたさまざまな2ブレーンとして、また大きさがゼロの $E_8\times E_8$ インスタントン上にも現れうる。一般に結合は1程度に固定されるので、摂動展開は存在しない。「張力がない」とは、重力スケールを単位とした弦の張力がゼロになることを意味する。この極限で相互作用しなくなる基本弦には、この用語を用いない。}

\gloss{テンソル多重項（tensor multiplet）}
{1. ボソン成分が一つの自己双対テンソルと一つのスカラーである、$d=6$、$(1,0)$ 超対称性の多重項。これは $d=4$、$N=2$ 超対称性のベクトル多重項に帰着する。2. ボソン成分が一つの自己双対テンソルと五つのスカラーである、$d=6$、$(2,0)$ 超対称性の多重項。これは $d=4$、$N=4$ 超対称性のベクトル多重項に帰着する。}

\gloss{テンソル演算子（tensor operator）}
{CFT において、共形変換または超共形変換に、変換の一階微分だけが現れる局所演算子。共形対称性または超共形対称性の一次場と同義である。このような演算子は、状態・演算子同型によって最高ウェイト状態へ写される。}

\gloss{テトラッド（tetrad）}
{$d$ 次元における、$d$ 個の正規直交ベクトル場の基底。tetrad という用語は、同義語の vierbein と同じく $d=4$ に由来するが、一般の $d$ に対する他の便利な用語はない。}

\gloss{かつて弦と呼ばれた理論（the theory formerly known as strings）}
{M 理論の第二の定義を参照。}

\gloss{理論（theory）}
{真空を参照。}

\gloss{テータ関数（theta functions）}
{トーラス上で単純な周期性をもつ正則関数の族。}

\gloss{Thirring 模型（Thirring model）}
{$1+1$ 次元で四次相互作用をもつ、単一の Dirac フェルミオンの可解な量子場の理論。ボソン化によって、一般の実数の半径をもつ自由スカラー場と等価になる。}

\gloss{しきい値補正（threshold correction）}
{仮想的な質量をもつ粒子による、低エネルギー有効作用、特にゲージ結合への補正。}

\gloss{位相的弦理論（topological string theory）}
{局所的な力学をもたないように修正された弦理論。すべての観測量は位相的である。}

\gloss{トーリック幾何（toric geometry）}
{射影空間の考え方の一般化。大まかにいえば、最も一般的な線形シグマ模型に対応する。}

\gloss{トーラスコンパクト化（toroidal compactification）}
{一つ以上の平坦な方向を周期的に同一視すること。}

\gloss{捩率（torsion）}
{さまざまな3形式の場の強さに用いられる用語。共変微分の中で Christoffel 接続と組み合わさって現れるため、この名で呼ばれる。}

\gloss{木レベル（tree-level）}
{一つの伝播関数を切ると非連結になる Feynman グラフ、またはそれに対応する弦振幅である球面と円板。これらは有効作用の古典的な項に対応する。}

\gloss{ツイスト（twist）}
{1. 動詞。周期的空間上の場を、対称変換 $h$ の違いだけ周期的でないものとして定義すること。2. 名詞。その非周期性 $h$。3. 動詞。CFT または弦理論が与えられたとき、対称群 $H$ を用いて新しい理論を構成すること。各元 $h\in H$ でツイストされた閉弦を加え、すべての状態が $H$ の変換に対して不変であることを要求する。これは $H$ を世界面のゲージ対称性として扱うことと等価である。オービフォールドという用語も同義語として用いられる。4. 名詞、古い用法。世界面パリティ。}

\gloss{ツイスト状態（twisted state）}
{ツイストされた周期性をもつ閉弦。}

\gloss{……型超重力（type \ldots\ supergravity）}
{対応する弦理論の低エネルギー超重力理論。ローマ数字は $d=10$ の超対称性の数を表し、IIA と IIB は、二つの超対称性のカイラリティがそれぞれ反対か同じかを区別する。}

\gloss{I 型超弦（type I superstring）}
{向きのない開超弦と閉超弦の理論であり、ゲージ群が $SO(32)$ の場合にのみ整合的である。右向きモードと左向きモードは、開弦の境界条件によって関係付けられるため、同じ時空超対称性のもとで変換する。}

\gloss{IIA 型超弦（type IIA superstring）}
{向きのある閉超弦の理論。右向きモードと左向きモードは、それぞれ別の時空超対称性のもとで変換し、それらのカイラリティは反対である。}

\gloss{IIB 型超弦（type IIB superstring）}
{向きのある閉超弦の理論。右向きモードと左向きモードは、それぞれ別の時空超対称性のもとで変換し、それらのカイラリティは同じである。}

\gloss{0 型弦（type 0 string）}
{タキオンをもち、時空フェルミオンをもたない、二つの非超対称な弦理論。II 型理論と世界面作用は同じだが、Hilbert 空間に対する射影が異なる。}

\gloss{U 双対性（U-duality）}
{弦理論、通常はトーラスコンパクト化した II 型理論における、あらゆる双対性。S 双対性と T 双対性を含むが、それらとは異なり、半径と結合を混ぜ合わせる変換も含む。}

\gloss{超短多重項（ultrashort multiplet）}
{BPS 状態を参照。}

\gloss{単位ゲージ（unit gauge）}
{2次元で、計量が単位行列になるように座標ゲージと Weyl ゲージを選ぶこと。}

\gloss{ユニタリ（unitary）}
{量子系について、正ノルムの Hilbert 空間に保存される内積があるという性質。}

\gloss{向きのない弦理論（unoriented string theory）}
{世界面のパリティ反転 $\Omega$ が離散的なゲージ対称性である弦理論。摂動論には向きのない世界面が含まれ、スペクトルは $\Omega=+1$ の状態に制限される。}

\gloss{紫外発散（UV divergence）}
{時空の短距離から生じる発散。通常は理論の有効性の限界を意味する。}

\gloss{真空（vacuum）}
{安定な Poincaré 不変状態。破れていない超対称性をもつ系の新しい特徴は、縮退しているが物理的に非等価な真空がしばしば現れることである。\par
異なる弦理論は、実際には単一の理論の異なる真空であることが今では明らかになっているが、それぞれに「理論」という語を用いるのが、なお一般的である。また、単一の弦理論の中の異なる CFT を、異なる真空ではなく異なる理論と呼ぶこともある。「模型」という語は、低エネルギーの物理が標準模型に似た弦の真空を指すために用いられる。}

\gloss{ベクトル多重項（vector multiplet）}
{$d=4$、$N=1$ または $d=6$、$N=1$ 超対称性において、ボソン成分が質量のないベクトル場である多重項。$d=6$ の多重項は、$d=4$、$N=2$ 超対称性では、ベクトル場と二つの実スカラーに帰着する。}

\gloss{Veneziano 振幅（Veneziano amplitude）}
{四つの開弦タキオンに対する、ボソン弦の木レベル振幅。}

\gloss{Verma 加群（Verma module）}
{最高ウェイト状態に、Virasoro の上昇生成子をあらゆる非等価な方法で作用させて得られる状態の集合。ただし、そのような状態がすべて線形独立であること、すなわち生成子が自由に作用することを要求する。これによって、中心電荷 $c$ と最高ウェイト状態のウェイト $h$ に依存する Virasoro 代数の表現が得られる。これは共形族と非常によく似ているが、縮退した表現の場合、共形族では一部の状態がゼロになりうる点が異なる。この場合、Verma 加群は可約表現を与える。}

\gloss{頂点演算子（vertex operator）}
{始状態または終状態の弦に対応する、弦の世界面上の局所演算子。}

\gloss{Virasoro 代数（Virasoro algebra）}
{CFT における、エネルギー運動量テンソルの Fourier モードの無限次元 Lie 代数。}

\gloss{Virasoro--Shapiro 振幅（Virasoro--Shapiro amplitude）}
{四つの閉弦タキオンに対する、ボソン弦の木レベル振幅。}

\gloss{$W$ 代数（$W$ algebras）}
{2より大きなスピンのカレントをもつ、拡張されたカイラル代数の族。}

\gloss{$W$ 弦（$W$ string）}
{2より大きなスピンの拘束をもつ、弦理論の一般化として提案されたもの。完全な構成は存在しない。}

\gloss{Ward 恒等式（Ward identity）}
{保存カレントを含む期待値の発散と、そのカレントを除いた同じ期待値との関係。}

\gloss{ウェイト（weight）}
{1. CFT において、スケール変換と回転のもとでの演算子の振る舞いを決める、$L_0$ または $\widetilde L_0$ の固有値。共形ウェイトと同じである。2. Lie 代数において、与えられた表現のある状態に対し、互いに可換な生成子の極大集合の固有値を並べたベクトル。}

\gloss{Wess--Zumino 整合性条件（Wess--Zumino consistency condition）}
{背景場に関する汎関数積分の二階微分が対称であるという条件。これは理論で可能なアノマリーの形を強く制約する。}

\gloss{Wess--Zumino--Novikov--Witten（WZNW）模型（Wess--Zumino--Novikov--Witten (WZNW) model）}
{反対称テンソル背景をもつ群多様体上の、共形不変な非線形シグマ模型。}

\gloss{Weyl アノマリー（Weyl anomaly）}
{Weyl 変換のアノマリー。Polyakov 形式では、これが臨界次元を決める。トレースアノマリー、または共形アノマリーと呼ぶこともある。}

\gloss{Weyl 条件（Weyl condition）}
{偶数次元で、フェルミオンが一定のカイラリティをもつという条件。これは Lorentz 群のスピノル表現を定義し、その成分数は、ガンマ行列代数の Dirac 表現の半分である。}

\gloss{Weyl 変換（Weyl transformation）}
{位置に依存する計量のスケール変換。}

\gloss{Wilson 線（Wilson line）}
{1. 場の強さはゼロだが、時空の非自明な経路を一周する平行移動、すなわちホロノミーが非自明であるゲージ場。2. そのような場を測るゲージ不変な演算子。ベクトルポテンシャルの線積分の、経路順序を付けた積のトレース。}

\gloss{Wilson 型作用（Wilsonian action）}
{より高いエネルギーの仮想状態の効果を取り込んだ、低エネルギー有効場の理論の作用。}

\gloss{巻き付き状態（winding state）}
{単連結でない時空において、その配置が非自明な経路となる閉弦。}

\gloss{世界面（world-sheet）}
{1. 弦の運動が時空内に掃く2次元曲面。2. 弦の運動を記述するために用いる、抽象的な2次元パラメーター空間。}

\gloss{世界面双対性（world-sheet duality）}
{弦の世界面を非等価な円に沿って切り開くことによって得られる、Hamiltonian による記述の間の等価性。重要な例として、球面の4点振幅から得られる結合律の条件や、円筒の開弦ループによる記述と閉弦の木による記述の等価性がある。}

\gloss{巻き付いた（wrapped）}
{$p$ ブレーンの $q$ 個の方向が、時空の非自明なコンパクト部分多様体に巻き付き、$(p-q)$ 次元の広がった対象が残ることをいう。}

\gloss{Yang--Mills 場（Yang--Mills field）}
{非 Abel ゲージ場。}

\gloss{Zamolodchikov 計量（Zamolodchikov metric）}
{球面上の一対の局所演算子の期待値であり、対応する弦状態の自然な内積を与える。これは時空作用の運動項に現れる計量である。}

\gloss{ゼロモード（zero modes）}
{1. 微分演算子のゼロモードとは、固有値がゼロの固有関数。Gauss 型汎関数積分では、ボソンの場合は無限大の因子、フェルミオンの場合はゼロの因子を与える。一般にこれらには物理的な起源があり、汎関数積分には有限の結果を与える適切な被積分関数がある。2. 重心モード。}

\gloss{零点エネルギー（zero-point energy）}
{量子場の真空揺らぎによるエネルギー。}

\gloss{ゼロスロープ極限（zero-slope limit）}
{$\alpha'\to0$ の極限。低エネルギー場の理論で記述される、質量のない弦状態だけが残る。}
